% Generated mechanically from frozen input inputs/p10/ja/spaces-resolve.tex.
% Do not edit this generated file; edit the bound locale source instead.

% OK, start here.
%

\setcounter{chapter}{88}
\chapter{曲面の特異点解消再訪}



\phantomsection
\label{spaces-resolve-section-phantom}


\section{序論}
\label{spaces-resolve-section-introduction}

\noindent
本章では，$2$ 次元 Noether 代数空間の特異点解消を論じる。
スキームに対する曲面の特異点解消については，既に前の章で
Lipman \cite{Lipman} に従って論じた。Resolution of Surfaces,
Section \ref{resolve-section-introduction} を参照せよ。
本章の結果の大部分は，スキームに対する結果から直ちに従う。

\medskip\noindent
特に断らない限り，本章の幾何学的対象はすべて代数空間とする。
したがって「$f : X \to Y$ を修正とする」と言うとき，これは
$f$ が Spaces over Fields, Definition
\ref{spaces-over-fields-definition-modification} の意味での射であることを表す。
proper 射などについても同様である。









\section{修正}
\label{spaces-resolve-section-modifications}

\noindent
$(A, \mathfrak m, \kappa)$ を Noether 局所環とする。
$S = \Spec(A)$ および $U = S \setminus \{\mathfrak m\}$ と置く。
本節では圏
\begin{equation}
\label{spaces-resolve-equation-modification}
\left\{
f : X \longrightarrow S
\quad \middle| \quad
\begin{matrix}
X\text{ は代数空間である}\\
f\text{ は proper 射である}\\
f^{-1}(U) \to U\text{ は同型である}
\end{matrix}
\right\}
\end{equation}
を考える。$X/S$ から $X'/S$ への射とは，$S$ 上の構造射と両立する
代数空間の射 $X \to X'$ とする。Algebraization of Formal Spaces,
Section \ref{restricted-section-modifications} では，この圏が $A$ の完備化
のみに依存することを見て，この圏の対象のいくつかの基本的性質を証明した。
本節では特に，$\dim(A) \leq 2$ の場合，または閉ファイバーの次元が
高々 $1$ の場合を調べる。

\begin{lemma}
\label{spaces-resolve-lemma-modification}
$(A, \mathfrak m, \kappa)$ を $2$ 次元 Noether 局所整域とし，
$U = \Spec(A) \setminus \{\mathfrak m\}$ は正規スキームであるとする。
このとき，任意の修正 $f : X \to \Spec(A)$ は
(\ref{spaces-resolve-equation-modification}) の射である。
\end{lemma}

\begin{proof}
$f : X \to S$ を修正とする。$f^{-1}(U) \to U$ が同型であることを
示せばよい。$U$ の任意の閉点 $u$ は余次元 $1$ なので，これは
Spaces over Fields, Lemma
\ref{spaces-over-fields-lemma-modification-normal-iso-over-codimension-1}
から従う。
\end{proof}

\begin{lemma}
\label{spaces-resolve-lemma-closed-immersion-on-fibre}
$(A, \mathfrak m, \kappa)$ を Noether 局所環とする。
$g : X \to Y$ を圏 (\ref{spaces-resolve-equation-modification}) の射とする。
特殊ファイバー上で誘導される射 $X_\kappa \to Y_\kappa$ が閉埋込みならば，
$g$ は閉埋込みである。
\end{lemma}

\begin{proof}
これは More on Morphisms of Spaces, Lemma
\ref{spaces-more-morphisms-lemma-where-closed-immersion} の特別な場合である。
\end{proof}

\begin{lemma}
\label{spaces-resolve-lemma-dimension-special-fibre}
$(A, \mathfrak m, \kappa)$ を次元 $\geq 1$ の Noether 局所整域とする。
$f : X \to \Spec(A)$ を代数空間の射とする。次の条件のうち少なくとも
一つが成り立つと仮定する：
\begin{enumerate}
\item $f$ は修正である（Spaces over Fields, Definition
\ref{spaces-over-fields-definition-modification}），
\item $f$ は alteration である（Spaces over Fields, Definition
\ref{spaces-over-fields-definition-alteration}），
\item $f$ は局所有限型かつ準分離的であり，$X$ は整で，
$\Spec(A)$ の生成点に写る $|X|$ の点がちょうど一つ存在する，
\item $f$ は局所有限型で，$X$ は decent であり，
$\Spec(A)$ の生成点に写る $|X|$ の点は $|X|$ の既約成分の生成点である，
% P10-E0403: frozen source contains the editorial placeholder below.
\item ここにさらに追加する。
\end{enumerate}
このとき $\dim(X_\kappa) \leq \dim(A) - 1$ である。
\end{lemma}

\begin{proof}
(1)，(2)，(3) は (4) の特別な場合である。アフィンスキーム
$U = \Spec(B)$ と \'etale 射 $U \to X$ を選ぶ。環準同型
$A \to B$ は有限型である。$\dim(U_\kappa) \leq \dim(A) - 1$
を示せばよい。$X$ は decent なので，$U$ の既約成分の生成点は
$|X|$ の既約成分の生成点の上にある点である。Decent Spaces, Lemma
\ref{decent-spaces-lemma-decent-generic-points} を参照せよ。
したがって $\Spec(B) \to \Spec(A)$ の $(0)$ 上のファイバーは，
$B$ の極小素イデアル $\mathfrak q_1, \ldots, \mathfrak q_r$ からなる
（有限）集合である。ゆえに，ある非零元 $f \in A$ に対して
$A_f \to B_f$ は有限である（Algebra, Lemma
\ref{algebra-lemma-generically-finite}）。従って
$\kappa(\mathfrak q_i)$ は $A$ の分数体の有限拡大である。
$\mathfrak m$ の上にある素イデアル $\mathfrak q \subset B$ を取る。このとき
$$
\dim(B_\mathfrak q) = \max \dim((B/\mathfrak q_i)_{\mathfrak q})
\leq \dim(A)
$$
である。不等式は $A \subset B/\mathfrak q_i$ に対する次元公式による。
Algebra, Lemma \ref{algebra-lemma-dimension-formula} を参照せよ。
しかし，$B_\mathfrak q/\mathfrak m B_\mathfrak q$ の次元
（これは対応する点での $U_\kappa$ の局所環の次元である）は，
極小素イデアル $\mathfrak q_i$ が $V(\mathfrak m)$ に含まれないため，
少なくとも一つ小さい。Properties, Lemma
\ref{properties-lemma-dimension} により結論を得る。
\end{proof}

\begin{lemma}
\label{spaces-resolve-lemma-modification-of-dim-2-is-projective-over-complete}
$(A, \mathfrak m, \kappa)$ が $2$ 次元完備 Noether 局所整域ならば，
$\Spec(A)$ の任意の修正は $A$ 上射影的である。
\end{lemma}

\begin{proof}
More on Morphisms of Spaces, Lemma
\ref{spaces-more-morphisms-lemma-projective-over-complete} により，
$\Spec(A)$ の任意の修正 $X$ の特殊ファイバーの次元が $\leq 1$
であることを示せば十分である。これは Lemma
\ref{spaces-resolve-lemma-dimension-special-fibre} から従う。
\end{proof}






\section{方針}
\label{spaces-resolve-section-strategy}

\noindent
$S$ をスキームとする。$X$ を $S$ 上の decent な代数空間とする。
$x_1, \ldots, x_n \in |X|$ を互いに異なる閉点とする。各 $i$ に対し，
Decent Spaces, Lemma
\ref{decent-spaces-lemma-decent-space-elementary-etale-neighbourhood}
の意味での初等 \'etale 近傍 $(U_i, u_i) \to (X, x_i)$ を選ぶ。
すなわち，$U_i$ はアフィンスキーム，$U_i \to X$ は \'etale，
$u_i$ は $x_i$ の上にある $U_i$ の唯一の点であり，
$\Spec(\kappa(u_i)) \to X$ は $x_i$ を表現するモノ射である。
$U_i$ を縮小することにより，$j \not = i$ のとき $x_j$ に写る
$U_i$ の点が存在しないと仮定してよいし，実際そう仮定する。
$u_i \in U_i$ は閉点であることに注意せよ。

\medskip\noindent
$f^{-1}(X \setminus \{x_1, \ldots, x_n\}) \to
X \setminus \{x_1, \ldots, x_n\}$ が同型となる代数空間の射
$f : Y \to X$ の圏を $\mathcal{C}_{X, \{x_1, \ldots, x_n\}}$
と書く。各 $i$ に対し，
$g_i^{-1}(U_i \setminus \{u_i\}) \to U_i \setminus \{u_i\}$
が同型となる代数空間の射 $g_i : Y_i \to U_i$ の圏を
$\mathcal{C}_{U_i, u_i}$ と書く。基底変換により関手
% P10-E0417: frozen source says "an functor"; the mathematical functor is retained.
\begin{equation}
\label{spaces-resolve-equation-equivalence}
F :
\mathcal{C}_{X, \{x_1, \ldots, x_n\}}
\longrightarrow
\mathcal{C}_{U_1, u_1} \times \ldots \times \mathcal{C}_{U_n, u_n}
\end{equation}
が定まる。本章の問題の少なくとも一部をスキームの場合に帰着するため，
次の補題を用いる。

\begin{lemma}
\label{spaces-resolve-lemma-equivalence}
関手 $F$ (\ref{spaces-resolve-equation-equivalence}) は圏同値である。
\end{lemma}

\begin{proof}
$n = 1$ の場合は Limits of Spaces, Lemma
\ref{spaces-limits-lemma-excision-modifications} である。
$n > 1$ の場合もまったく同じ方法で証明でき，また前者から
導くこともできる。例えば，$g_i : Y_i \to U_i$ が
$\mathcal{C}_{U_i, u_i}$ の対象であるとする。$n = 1$ の場合により，
$X \setminus \{x_i\}$ 上で同型で，$U_i$ への基底変換が
% P10-E0404: frozen source names the undefined f_i rather than the given g_i.
$f_i$ となる $f'_i : Y'_i \to X$ を取れる。このとき
$$
f : Y = Y'_1 \times_X \ldots \times_X Y'_n \to X
$$
と置ける。これは $\mathcal{C}_{X, \{x_1, \ldots, x_n\}}$ の対象で，
$U_i \to X$ による基底変換は $g_i$ を復元する。したがって関手は
本質的全射である。充満忠実性の証明は省略する。
% P10-E0405: frozen English has the malformed phrase "fully faithfulness".
\end{proof}

\begin{lemma}
\label{spaces-resolve-lemma-equivalence-properties}
$X, x_i, U_i \to X, u_i$ を (\ref{spaces-resolve-equation-equivalence}) の通りとする。
$f : Y \to X$ が $F$ の下で $g_i : Y_i \to U_i$ に対応するとする。
このとき，$f$ が準コンパクト，準分離的，分離的，局所有限表示，有限表示，
局所有限型，有限型，proper，整，有限のいずれかであることと，
各 $i = 1, \ldots, n$ について $g_i$ がその性質をもつことは同値である。
\end{lemma}

\begin{proof}
Limits of Spaces, Lemma
\ref{spaces-limits-lemma-excision-modifications-properties} から従う。
\end{proof}

\begin{lemma}
\label{spaces-resolve-lemma-equivalence-fibre}
$X, x_i, U_i \to X, u_i$ を (\ref{spaces-resolve-equation-equivalence}) の通りとする。
$f : Y \to X$ が $F$ の下で $g_i : Y_i \to U_i$ に対応するとする。
このとき代数空間として $Y_{x_i} \cong (Y_i)_{u_i}$ である。
\end{lemma}

\begin{proof}
$u_i \to x_i$ は同型なので明らかである。
\end{proof}






\section{二次変換による支配}
\label{spaces-resolve-section-quadratic-spaces}

\noindent
空間の点におけるブローアップは，$X$ が decent の場合にのみ定義する。

\begin{definition}
\label{spaces-resolve-definition-blowup-at-point}
$S$ をスキームとし，$X$ を $S$ 上の decent な代数空間とする。
$x \in |X|$ を閉点とする。Decent Spaces, Lemma
\ref{decent-spaces-lemma-decent-space-closed-point} により，
$x$ は閉埋込み $i : \Spec(k) \to X$ で表現できる。
閉部分空間 $Z = i(\Spec(k)) \subset X$ における $X$ のブローアップを，
{\it $x$ における $X$ のブローアップ $X' \to X$} と呼ぶ。
\end{definition}

\noindent
この一般性では，$x$ における $X$ のブローアップは必ずしも proper ではない。
しかし $X$ が局所 Noether ならば，Divisors on Spaces, Lemma
\ref{spaces-divisors-lemma-blowing-up-projective} によりブローアップは proper である。
局所 Noether 代数空間が Noether であるための必要十分条件は，
それが準コンパクトかつ準分離的であることであった。また局所 Noether
代数空間について，準分離的であることと decent であることは同値である
（Decent Spaces, Lemma
\ref{decent-spaces-lemma-locally-Noetherian-decent-quasi-separated}）。

\begin{lemma}
\label{spaces-resolve-lemma-equivalence-sequence-blowups}
$X, x_i, U_i \to X, u_i$ を (\ref{spaces-resolve-equation-equivalence}) の通りとし，
$f : Y \to X$ が $F$ の下で $g_i : Y_i \to U_i$ に対応すると仮定する。
このとき，$f$ が分解
$$
Y = Z_m \to Z_{m - 1} \to \ldots \to Z_1 \to Z_0 = X
$$
をもち，各 $Z_{j + 1} \to Z_j$ が
$\{x_1, \ldots, x_n\}$ の上にある閉点 $z_j$ における $Z_j$ の
ブローアップであることと，各 $i$ について $g_i$ が分解
$$
Y_i = Z_{i, m_i} \to Z_{i, m_i - 1} \to \ldots \to Z_{i, 1} \to Z_{i, 0} = U_i
$$
をもち，各 $Z_{i, j + 1} \to Z_{i, j}$ が $u_i$ の上にある閉点
$z_{i, j}$ における $Z_{i, j}$ のブローアップであることは同値である。
\end{lemma}

\begin{proof}
ブローアップは表現可能な射である。したがっていずれの場合も帰納的に，
$Z_j \to X$ または $Z_{i, j} \to U_i$ が表現可能であることが分かる。
ゆえに Decent Spaces, Lemma
\ref{decent-spaces-lemma-representable-named-properties} により，
各 $Z_j$ または $Z_{i, j}$ は decent な代数空間である。
これにより主張は意味をもつ（ブローアップは decent な空間に対してのみ
定義されているからである）。同値性を証明するため，まずブローアップの列
$Z_m \to Z_{m - 1} \to \ldots \to Z_1 \to Z_0 = X$ から始める。
最初の射 $Z_1 \to X$ は，ある $x_i$，例えば $x_1$ のブローアップで与えられる。
$F$ を $Z_1 \to X$ に適用すると，
% P10-E0406: the frozen source repeats the predicate in this sentence.
$u_1$ におけるブローアップ $Z_{1, 1} \to U_1$ が $u_1$ における
ブローアップであり，それ以外では $i > 1$ に対して
$Z_{i, 0} = U_i$ であることが分かる。次の段階では，$Z_1$ 上で
$x_i$ の一つ（$i \geq 2$）をブローアップするか，または
$x_1$ 上の $Z_1 \to X$ のファイバーの閉点 $z_1$ を選ぶ。
前者では何をすべきかは明らかであり，後者では
$(Z_1)_{x_1} \cong (Z_{1, 1})_{u_1}$
（Lemma \ref{spaces-resolve-lemma-equivalence-fibre}）を用いて，$z_1$ に対応する閉点
$z_{1, 1} \in Z_{1, 1}$ を得る。次に $Z_{1, 2} \to Z_{1, 1}$ を
$z_{1, 1}$ におけるブローアップとする。このように続けて，各 $g_i$ の
分解を構成する。

\medskip\noindent
逆に，ブローアップの列
$Z_{i, m_i} \to Z_{i, m_i - 1} \to \ldots \to Z_{i, 1} \to Z_{i, 0} = U_i$
が与えられたときも，まったく同じ方法で $X$ のブローアップの列を構成する。
\end{proof}

\begin{lemma}
\label{spaces-resolve-lemma-make-ideal-principal}
$S$ をスキームとし，$X$ を $S$ 上の Noether 代数空間とする。
$T \subset |X|$ を，(1) $X$ が $x$ で正則であり，かつ
(2) $x$ における $X$ の局所環の次元が $2$ であるような閉点 $x$ の
有限集合とする。$\mathcal{I} \subset \mathcal{O}_X$ を，
$\mathcal{O}_X/\mathcal{I}$ が $T$ 上に台をもつ準連接イデアル層とする。
このとき列
$$
X_m \to X_{m - 1} \to \ldots \to X_1 \to X_0 = X
$$
であって，各 $X_{j + 1} \to X_j$ が $T$ の点の上にある閉点 $x_j$ に
おける $X_j$ のブローアップであり，
% P10-E0407: the frozen source switches from the displayed X_m to X_n.
$\mathcal{I}\mathcal{O}_{X_n}$ が可逆イデアル層となるものが存在する。
\end{lemma}

\begin{proof}
$T = \{x_1, \ldots, x_r\}$ とする。Section \ref{spaces-resolve-section-strategy}
の通り，初等 \'etale 近傍
$(U_i, u_i) \to (X, x_i)$ を選ぶ。各 $i$ について制限
$\mathcal{I}_i = \mathcal{I}|_{U_i} \subset \mathcal{O}_{U_i}$ は
$u_i$ に台をもつ準連接イデアル層である。$u_i$ における $U_i$ の局所環は
正則で次元 $2$ である。したがって Resolution of Surfaces, Lemma
\ref{resolve-lemma-make-ideal-principal} を適用して，
% P10-E0408: the frozen i-indexed sequence uses the unindexed X_1.
$$
X_{i, m_i} \to X_{i, m_i - 1} \to \ldots \to X_1 \to X_{i, 0} = U_i
$$
という，$u_i$ の上の閉点におけるブローアップの列で，
$\mathcal{I}_i \mathcal{O}_{X_{i, m_i}}$ が可逆となるものを得る。
Lemma \ref{spaces-resolve-lemma-equivalence-sequence-blowups} により，主張の通りの
ブローアップの列
$$
X_m \to X_{m - 1} \to \ldots \to X_1 \to X_0 = X
$$
で，$U_i$ への基底変換が与えられた列を生じるものを得る。
このとき $\mathcal{I}\mathcal{O}_{X_n}$ は可逆イデアル層である。
すなわち，これは
% P10-E0409: the frozen source uses the undefined n instead of r here.
$X \setminus \{x_1, \ldots, x_n\}$ 上で成り立ち，また各 $x_i$ の
ファイバーの \'etale 近傍では構成により成り立つ。
\end{proof}

\begin{lemma}
\label{spaces-resolve-lemma-dominate-by-blowing-up-in-points}
$S$ をスキームとし，$X$ を $S$ 上の Noether 代数空間とする。
$T \subset |X|$ を，(1) $X$ が $x$ で正則であり，かつ
(2) $x$ における $X$ の局所環の次元が $2$ であるような閉点 $x$ の
有限集合とする。$f : Y \to X$ を，$U = X \setminus T$ 上で同型となる
代数空間の proper 射とする。このとき列
$$
X_n \to X_{n - 1} \to \ldots \to X_1 \to X_0 = X
$$
であって，各 $X_{i + 1} \to X_i$ が $T$ の点の上にある閉点 $x_i$ に
おける $X_i$ のブローアップであり，合成が $X_n \to Y \to X$ と
分解するものが存在する。
\end{lemma}

\begin{proof}
More on Morphisms of Spaces, Lemma
\ref{spaces-more-morphisms-lemma-dominate-modification-by-blowup}
により，$Y \to X$ を支配する $U$-許容ブローアップ $X' \to X$ が存在する。
したがって，$\mathcal{O}_X/\mathcal{I}$ が $T$ 上に台をもち，
$Y$ が $\mathcal{I}$ における $X$ のブローアップとなるような
イデアル層 $\mathcal{I} \subset \mathcal{O}_X$ が存在すると仮定してよい。
Lemma \ref{spaces-resolve-lemma-make-ideal-principal} により列
$$
X_n \to X_{n - 1} \to \ldots \to X_1 \to X_0 = X
$$
で，各 $X_{i + 1} \to X_i$ が $T$ の点の上にある閉点 $x_i$ における
$X_i$ のブローアップであり，$\mathcal{I}\mathcal{O}_{X_n}$ が
可逆イデアル層となるものが存在する。ブローアップの普遍性
（Divisors on Spaces, Lemma
\ref{spaces-divisors-lemma-universal-property-blowing-up}）により，
所望の分解を得る。
\end{proof}





\section{正規化ブローアップによる支配}
\label{spaces-resolve-section-normalized-blowups}

\noindent
本節では，曲面の修正が点における正規化ブローアップの列によって
支配されることを証明する。

\begin{definition}
\label{spaces-resolve-definition-normalized-blowup}
$S$ をスキームとする。$X$ を，Morphisms of Spaces, Lemma
\ref{spaces-morphisms-lemma-prepare-normalization} の同値な条件を満たす
$S$ 上の decent な代数空間とする。$x \in |X|$ を閉点とする。
合成 $X'' \to X' \to X$ を {\it $x$ における $X$ の正規化ブローアップ}
と呼ぶ。ただし $X' \to X$ は $x$ における $X$ のブローアップ
（Definition \ref{spaces-resolve-definition-blowup-at-point}）であり，
$X'' \to X'$ は $X'$ の正規化である。
\end{definition}

\noindent
% P10-E0410: the frozen source leaves the causal clause grammatically unattached.
ここで正規化 $X'' \to X'$ は代数空間として定義される。$X'$ が
Morphisms of Spaces, Lemma
\ref{spaces-morphisms-lemma-prepare-normalization} の同値な条件を満たすことは，
Divisors on Spaces, Lemma
\ref{spaces-divisors-lemma-blowup-finite-nr-irreducibles} による。
正規化の定義については Morphisms of Spaces, Definition
\ref{spaces-morphisms-definition-normalization} を参照せよ。

\medskip\noindent
一般に，$X$ が Noether であっても正規化ブローアップは proper とは限らない。
代数空間が Nagata であるとは，Nagata 環のスペクトルであるアフィンによる
\'etale 被覆をもつことであった（Properties of Spaces, Definition
\ref{spaces-properties-definition-type-property}，Remark
\ref{spaces-properties-remark-list-properties-local-etale-topology}，および
Properties, Definition \ref{properties-definition-nagata}）。

\begin{lemma}
\label{spaces-resolve-lemma-Nagata-normalized-blowup}
Definition \ref{spaces-resolve-definition-normalized-blowup} において $X$ が Nagata ならば，
$x$ における $X$ の正規化ブローアップは，$X$ 上 proper な正規 Nagata
代数空間である。
\end{lemma}

\begin{proof}
ブローアップ射 $X' \to X$ は proper である（$X$ は局所 Noether なので
Divisors on Spaces, Lemma
\ref{spaces-divisors-lemma-blowing-up-projective} を適用できる）。
したがって $X'$ は Nagata である（Morphisms of Spaces, Lemma
\ref{spaces-morphisms-lemma-finite-type-nagata}）。ゆえに正規化
$X'' \to X'$ は有限であり（Morphisms of Spaces, Lemma
\ref{spaces-morphisms-lemma-nagata-normalization}），従って
$X'' \to X$ も proper である（Morphisms of Spaces, Lemmas
\ref{spaces-morphisms-lemma-finite-proper} および
\ref{spaces-morphisms-lemma-composition-proper}）。正規化ブローアップは
正規（Morphisms of Spaces, Lemma
\ref{spaces-morphisms-lemma-normalization-normal}）Nagata 代数空間である。
\end{proof}

\noindent
次は Lemma \ref{spaces-resolve-lemma-equivalence-sequence-blowups} の正規化ブローアップ版である。

\begin{lemma}
\label{spaces-resolve-lemma-equivalence-sequence-normalized-blowups}
$X, x_i, U_i \to X, u_i$ を (\ref{spaces-resolve-equation-equivalence}) の通りとし，
$f : Y \to X$ が $F$ の下で $g_i : Y_i \to U_i$ に対応すると仮定する。
$X$ は Morphisms of Spaces, Lemma
\ref{spaces-morphisms-lemma-prepare-normalization} の同値な条件を満たすとする。
このとき，$f$ が分解
$$
Y = Z_m \to Z_{m - 1} \to \ldots \to Z_1 \to Z_0 = X
$$
をもち，各 $Z_{j + 1} \to Z_j$ が
$\{x_1, \ldots, x_n\}$ の上にある閉点 $z_j$ における $Z_j$ の
正規化ブローアップであることと，各 $i$ について $g_i$ が分解
$$
Y_i = Z_{i, m_i} \to Z_{i, m_i - 1} \to \ldots \to Z_{i, 1} \to Z_{i, 0} = U_i
$$
をもち，各 $Z_{i, j + 1} \to Z_{i, j}$ が $u_i$ の上にある閉点
$z_{i, j}$ における $Z_{i, j}$ の正規化ブローアップであることは同値である。
\end{lemma}

\begin{proof}
Lemma \ref{spaces-resolve-lemma-equivalence-sequence-blowups} の証明とまったく同じ議論で従う。
\end{proof}

\noindent
Nagata 代数空間は局所 Noether である。

\begin{lemma}
\label{spaces-resolve-lemma-dominate-by-normalized-blowing-up}
$S$ をスキームとする。$X$ を $S$ 上で $\dim(X) = 2$ を満たす Noether Nagata
代数空間とする。$f : Y \to X$ を proper 双有理射とする。
このとき可換図式
$$
\xymatrix{
X_n \ar[r] \ar[d] &
X_{n - 1} \ar[r] &
\ldots \ar[r] &
X_1 \ar[r] &
X_0 \ar[d] \\
Y \ar[rrrr]  & & & & X
}
$$
であって，$X_0 \to X$ が正規化であり，各 $X_{i + 1} \to X_i$ が
$X_i$ の閉点における正規化ブローアップであるものが存在する。
\end{lemma}

\begin{proof}
この補題は代数空間について直接証明することもできるが，上で用いた方法を
続けてスキームの場合に帰着する。

\medskip\noindent
Noether 代数空間は準分離的なので点の剰余体がよく定義されることを用いる
（例えば Decent Spaces, Lemma
\ref{decent-spaces-lemma-decent-space-elementary-etale-neighbourhood} による）。
以下では Morphisms of Spaces, Sections
\ref{spaces-morphisms-section-nagata},
\ref{spaces-morphisms-section-dimension-formula}, and
\ref{spaces-morphisms-section-normalization} の結果を断りなく用いる。
$Y$ をその正規化で置き換えてよい。$X_0 \to X$ を正規化とする。
射 $Y \to X$ は $X_0$ を経由する。したがって $X$ と $Y$ はともに
正規であると仮定してよい。

\medskip\noindent
$X$ と $Y$ が正規であると仮定する。射 $f : Y \to X$ は，$Y$ の
余次元 $0$ および $1$ のすべての点と，その上のファイバーが有限である
$Y$ のすべての点を含む開集合上で同型である。Spaces over Fields, Lemma
\ref{spaces-over-fields-lemma-modification-normal-iso-over-codimension-1}
を参照せよ。従って，$f$ が $X \setminus T$ 上で同型となるような
閉点の有限集合 $T \subset |X|$ が存在する。More on Morphisms of Spaces,
Lemma \ref{spaces-more-morphisms-lemma-dominate-modification-by-blowup}
により，$Y$ を支配する $X \setminus T$-許容ブローアップ
$Y' \to X$ が存在する。$Y$ を $Y'$ の正規化で置き換えると，
$Y \to X$ は表現可能であると仮定してよい。

\medskip\noindent
$T = \{x_1, \ldots, x_r\}$ とする。Section \ref{spaces-resolve-section-strategy}
の通り，初等 \'etale 近傍 $(U_i, u_i) \to (X, x_i)$ を選ぶ。
各 $i$ について，射 $Y_i = Y \times_X U_i \to U_i$ は
$U_i \setminus \{u_i\}$ 上で同型となる proper 双有理射である。
したがって Resolution of Surfaces, Lemma
\ref{resolve-lemma-dominate-by-normalized-blowing-up} を適用して，
% P10-E0411: the frozen i-indexed sequence uses the unindexed X_1.
$$
X_{i, m_i} \to X_{i, m_i - 1} \to \ldots \to X_1 \to X_{i, 0} = U_i
$$
という，$u_i$ の上の閉点における正規化ブローアップの列で，
$X_{i, m_i}$ が $Y_i$ を支配するものを得る。
Lemma \ref{spaces-resolve-lemma-equivalence-sequence-normalized-blowups} により，
主張の通りの正規化ブローアップの列
$$
X_m \to X_{m - 1} \to \ldots \to X_1 \to X_0 = X
$$
で，$U_i$ への基底変換が与えられた列を生じるものを得る。
Lemma \ref{spaces-resolve-lemma-equivalence} の圏同値により，$X_m$ は $Y$ を支配する。
\end{proof}








\section{完備化への基底変換}
\label{spaces-resolve-section-aux}

\noindent
次の簡単な補題は，以下で有用な道具となる。

\begin{lemma}
\label{spaces-resolve-lemma-iso-completions}
$(A, \mathfrak m, \kappa)$ を，極大イデアル $\mathfrak m$ が有限生成である
局所環とする。$X$ を $A$ 上の decent な代数空間とする。
$A^\wedge$ を $A$ の $\mathfrak m$-進完備化とし，
$Y = X \times_{\Spec(A)} \Spec(A^\wedge)$ と置く。
$q \in |Y|$ の像 $p \in |X|$ が $\Spec(A)$ の閉点の上にあるとする。
このとき Hensel 局所環の準同型
$\mathcal{O}_{X, p}^h \to \mathcal{O}_{Y, q}^h$ は，完備化上の同型を誘導する。
\end{lemma}

\begin{proof}
Decent Spaces, Lemma
\ref{decent-spaces-lemma-decent-space-elementary-etale-neighbourhood}
の通り初等 \'etale 近傍 $(U, u) \to (X, p)$ を選び，
$V = U \times_X Y = U \times_{\Spec(A)} \Spec(A^\wedge)$ および
$v = (u, q)$ と置けば，スキームの場合から直ちに従う。
スキームの場合は Resolution of Surfaces, Lemma
\ref{resolve-lemma-iso-completions} である。
\end{proof}

\begin{lemma}
\label{spaces-resolve-lemma-port-regularity-to-completion}
$(A, \mathfrak m, \kappa)$ を Noether 局所環とする。
$X$ を decent な代数空間とし，$X \to \Spec(A)$ を局所有限型の射とする。
$Y = X \times_{\Spec(A)} \Spec(A^\wedge)$ と置く。
$y \in |Y|$ の像を $x \in |X|$ とする。このとき
\begin{enumerate}
\item $\mathcal{O}_{Y, y}^h$ が正則ならば，$\mathcal{O}_{X, x}^h$ は正則である，
\item $y$ が閉ファイバーに属するならば，$\mathcal{O}_{Y, y}^h$ が正則
$\Leftrightarrow \mathcal{O}_{X, x}^h$ が正則である，
\item $X$ が $A$ 上 proper ならば，$X$ が正則であることと $Y$ が正則であることは同値である。
\end{enumerate}
\end{lemma}

\begin{proof}
\'etale 局所化により，最初の二つの主張はスキームに対する対応する補題から
直ちに従う。Resolution of Surfaces, Lemma
\ref{resolve-lemma-port-regularity-to-completion} を参照せよ。
% P10-E0418: frozen English splits the compound noun "counterpart".
(3) について，$A \to A^\wedge$ は忠実平坦なので $Y \to X$ は全射であり，
(1) により $Y$ が正則ならば $X$ は正則である。逆に $X$ が正則ならば，
特殊ファイバーのすべての点において $Y$ の Hensel 局所環は正則である。
一般の点 $y \in |Y|$ を取る。proper の場合
$|Y| \to |\Spec(A^\wedge)|$ は閉写像なので，閉ファイバーにある
$y_0$ への特殊化 $y \leadsto y_0$ を取れる。Decent Spaces, Lemma
\ref{decent-spaces-lemma-decent-space-elementary-etale-neighbourhood}
の通り初等 \'etale 近傍 $(V, v_0) \to (Y, y_0)$ を選ぶ。
$Y$ は decent なので，$y \leadsto y_0$ を $V$ における特殊化
$v \leadsto v_0$ に持ち上げられる（Decent Spaces, Lemma
\ref{decent-spaces-lemma-decent-specialization}）。従って
$\mathcal{O}_{V, v}$ は $\mathcal{O}_{V, v_0}$ の局所化であり，ゆえに正則である。
これで証明が完了する。
\end{proof}

\begin{lemma}
\label{spaces-resolve-lemma-formally-unramified}
$(A, \mathfrak m)$ を Noether 局所環とし，$X$ を $A$ 上の代数空間とする。
次を仮定する：
\begin{enumerate}
\item $A$ は解析的不分岐である
（Algebra, Definition \ref{algebra-definition-analytically-unramified}），
\item $X$ は $A$ 上局所有限型である，
\item $X \to \Spec(A)$ は $X$ のすべての余次元 $0$ の点で \'etale である。
\end{enumerate}
このとき $X$ の正規化は $X$ 上有限である。
\end{lemma}

\begin{proof}
スキーム $U$ と全射 \'etale 射 $U \to X$ を選ぶ。このとき
$U \to \Spec(A)$ は Resolution of Surfaces, Lemma
\ref{resolve-lemma-formally-unramified} の仮定を満たすので，その結論も満たす。
\end{proof}










\section{存在から従う性質}
\label{spaces-resolve-section-existence-gives}

\noindent
本節では，Noether 整代数空間について，正則 alteration の存在から
多くの帰結が得られることを証明する。「悪い」Noether 代数空間に
関心のない読者は本節を飛ばしてよい。

\begin{lemma}
\label{spaces-resolve-lemma-regular-alteration-implies}
$S$ をスキームとし，$Y$ を $S$ 上の Noether 整代数空間とする。
$X$ が正則であるような alteration $f : X \to Y$ が存在すると仮定する。
このとき正規化 $Y^\nu \to Y$ は有限であり，$Y$ は正則な稠密開集合をもつ。
\end{lemma}

\begin{proof}
\'etale 局所化により，$Y = \Spec(A)$，$A$ が Noether 整域の場合を
証明すれば十分である。$B$ を $A$ の分数体における整閉包とし，
$C = \Gamma(X, \mathcal{O}_X)$ と置く。Cohomology of Spaces, Lemma
\ref{spaces-cohomology-lemma-proper-pushforward-coherent} により，
$C$ は有限 $A$-加群である。$X$ は正規なので（Properties of Spaces,
Lemma \ref{spaces-properties-lemma-regular-normal}），
% P10-E0419: frozen English omits the article before "normal domain" here.
$C$ は正規整域である（Spaces over Fields, Lemma
\ref{spaces-over-fields-lemma-normal-integral-sections}）。
従って $B \subset C$ であり，$A$ は Noether なので $B$ は $A$ 上有限である。

\medskip\noindent
$f^{-1}V \to V$ が有限となる空でない開集合 $V \subset Y$ が存在する。
Spaces over Fields, Definition
\ref{spaces-over-fields-definition-alteration} を参照せよ。
$V$ を縮小して，$f^{-1}V \to V$ は平坦であると仮定してよい
（Morphisms of Spaces, Proposition
\ref{spaces-morphisms-proposition-generic-flatness-reduced}）。
従って $f^{-1}V \to V$ は忠実平坦である。Algebra, Lemma
\ref{algebra-lemma-descent-regular} により $V$ は正則である。
\end{proof}

\begin{lemma}
\label{spaces-resolve-lemma-regular-alteration-implies-local}
$(A, \mathfrak m, \kappa)$ を Noether 局所整域とする。
$X$ が正則であるような alteration $f : X \to \Spec(A)$ が存在すると仮定する。
このとき
\begin{enumerate}
\item ある非零元 $f \in A$ に対して $A_f$ は正則である，
\item $A$ の分数体における整閉包 $B$ は $A$ 上有限である，
\item $B$ の $\mathfrak m$-進完備化は正規環である。すなわち，
$B$ の極大イデアルにおける完備化は正規整域である，
\item $A$ の生成形式ファイバーは正則である。
\end{enumerate}
\end{lemma}

\begin{proof}
(1) と (2) は Lemma \ref{spaces-resolve-lemma-regular-alteration-implies} から従う。
(3) の証明の記法を準備するため，その補題の証明の一部を繰り返す必要がある。
$C = \Gamma(X, \mathcal{O}_X)$ と置く。Cohomology of Spaces, Lemma
\ref{spaces-cohomology-lemma-proper-pushforward-coherent} により，
$C$ は有限 $A$-加群である。$X$ は正規なので（Properties of Spaces,
Lemma \ref{spaces-properties-lemma-regular-normal}），
% P10-E0419: the second frozen occurrence also omits the article.
$C$ は正規整域である（Spaces over Fields, Lemma
\ref{spaces-over-fields-lemma-normal-integral-sections}）。
従って $B \subset C$ であり，$A$ は Noether なので $B$ は $A$ 上有限である。
Resolution of Surfaces, Lemma \ref{resolve-lemma-algebra-helper} により，
(3) を証明するには $\mathfrak m$-進完備化 $C^\wedge$ が正規であることを
示せば十分である。

\medskip\noindent
Algebra, Lemma \ref{algebra-lemma-completion-finite-extension} により，
完備化 $C^\wedge$ は，$\mathfrak m$ の上にある $C$ の素イデアルにおける
$C$ の完備化の積である。このような素イデアルは有限個で，$C$ の極大イデアル
$\mathfrak m_1, \ldots, \mathfrak m_r$ である。
（$B$ に対する対応する結果が補題の最後の主張を説明する。）
従って $A$ を $C_{\mathfrak m_i}$ で，$X$ を
$X_i = X \times_{\Spec(C)} \Spec(C_{\mathfrak m_i})$ で置き換えることにより，
次段落の場合に帰着する。（Cohomology of Spaces, Lemma
\ref{spaces-cohomology-lemma-flat-base-change-cohomology} により
$\Gamma(X_i, \mathcal{O}) = C_{\mathfrak m_i}$ であることに注意せよ。）

\medskip\noindent
ここで $A$ は Noether 局所正規整域であり，$f : X \to \Spec(A)$ は
$\Gamma(X, \mathcal{O}_X) = A$ を満たす正則 alteration である。
$A$ の完備化 $A^\wedge$ が正規整域であることを示さなければならない。
Lemma \ref{spaces-resolve-lemma-port-regularity-to-completion} により
$Y = X \times_{\Spec(A)} \Spec(A^\wedge)$ は正則である。
% P10-E0412: the frozen source sentence beginning with Since has no main clause.
Cohomology of Spaces, Lemma
\ref{spaces-cohomology-lemma-flat-base-change-cohomology} により
$\Gamma(Y, \mathcal{O}_Y) = A^\wedge$ であるので。
前と同様に $A^\wedge$ は正規であると結論する。実際，Properties of Spaces,
Lemma \ref{spaces-properties-lemma-regular-normal} により $Y$ は正規である。
$\Gamma(Y, \mathcal{O}_Y) = A^\wedge$ は局所環なので，この空間は連結である。
ゆえに $Y$ は正規かつ整である（分離的代数空間について，連結かつ正規ならば
整である）。従って Spaces over Fields, Lemma
\ref{spaces-over-fields-lemma-normal-integral-sections} により
$\Gamma(Y, \mathcal{O}_Y) = A^\wedge$ は正規整域である。これで (3) を証明した。

\medskip\noindent
(4) の証明。$\eta \in \Spec(A)$ を生成点とし，下付き添字 $\eta$ で
$\eta$ への基底変換を表す。$f$ は alteration なので，スキーム
$X_\eta$ は $\eta$ 上有限忠実平坦である。Lemma
\ref{spaces-resolve-lemma-port-regularity-to-completion} により
$Y = X \times_{\Spec(A)} \Spec(A^\wedge)$ は正則なので，
$Y_\eta$ は正則である（$Y$ の開集合の極限として）。このとき
$Y_\eta \to \Spec(A^\wedge \otimes_A \kappa(\eta))$ は生成形式ファイバー上
有限忠実平坦である。Algebra, Lemma
\ref{algebra-lemma-descent-regular} により結論を得る。
\end{proof}










\section{特異点解消}
\label{spaces-resolve-section-resolution}

\noindent
まず定義を与える。

\begin{definition}
\label{spaces-resolve-definition-resolution}
$S$ をスキームとし，$Y$ を $S$ 上の Noether 整代数空間とする。
% P10-E0413: frozen source introduces Y but calls this a resolution of X.
$X$ の {\it 特異点解消} とは，$X$ が正則であるような修正
$f : X \to Y$ のことである。
\end{definition}

\noindent
曲面の場合には，もう少し詳しい情報が必要となることがある。

\begin{definition}
\label{spaces-resolve-definition-resolution-surface}
$S$ をスキームとし，$Y$ を $S$ 上の $2$ 次元 Noether 整代数空間とする。
$Y$ が {\it 正規化ブローアップによる特異点解消} をもつとは，列
% P10-E0414: frozen sequence mixes Y_i and X_{n-1}; preserve exact symbols.
$$
Y_n \to X_{n - 1} \to \ldots \to Y_1 \to Y_0 \to Y
$$
が存在し，次を満たすことをいう：
\begin{enumerate}
\item $i = 0, \ldots, n$ に対して $Y_i$ は $Y$ 上 proper である，
\item $Y_0 \to Y$ は正規化である，
\item $i = 1, \ldots, n$ に対して $Y_i \to Y_{i - 1}$ は正規化ブローアップである，
\item $Y_n$ は正則である。
\end{enumerate}
\end{definition}

\noindent
条件 (1) から，$Y$ の正規化 $Y_0$ は $Y$ 上有限であり，また正規化
ブローアップで用いる正規化も有限であることに注意せよ。
これで本章の主定理に到達した。

\begin{theorem}
\label{spaces-resolve-theorem-resolve}
$S$ をスキームとし，$Y$ を $S$ 上の 2 次元整 Noether 代数空間とする。
次は同値である：
\begin{enumerate}
\item $X$ が正則であるような alteration $X \to Y$ が存在する，
\item $Y$ の特異点解消が存在する，
\item $Y$ は正規化ブローアップによる特異点解消をもつ，
\item 正規化 $Y^\nu \to Y$ は有限であり，$Y^\nu$ は有限個の特異点
% P10-E0415: frozen source locates points of Y^nu in |Y|.
$y_1, \ldots, y_m \in |Y|$ をもち，各 $i$ について Hensel 局所環
$\mathcal{O}_{Y^\nu, y_i}^h$ の完備化は正規である。
\end{enumerate}
\end{theorem}

\begin{proof}
含意 (3) $\Rightarrow$ (2) $\Rightarrow$ (1) は明らかである。

\medskip\noindent
$X$ が正則であるような alteration $X \to Y$ を取る。Lemma
\ref{spaces-resolve-lemma-regular-alteration-implies} により $Y^\nu \to Y$ は有限である。
Morphisms of Spaces, Lemma
\ref{spaces-morphisms-lemma-normalization-normal} から得られる分解
$f : X \to Y^\nu$ を考える。Spaces over Fields, Lemma
\ref{spaces-over-fields-lemma-finite-in-codim-1} により，$f$ は
$Y^\nu$ のすべての余次元 $\leq 1$ の点を含む開集合
$V \subset Y^\nu$ 上で有限である。Algebra, Lemma
\ref{algebra-lemma-CM-over-regular-flat} と，Serre の判定法により次元
$\leq 2$ の正規局所環が Cohen--Macaulay であること（Algebra, Lemma
\ref{algebra-lemma-criterion-normal}）から，$f$ は $V$ 上平坦である。
Algebra, Lemma \ref{algebra-lemma-descent-regular} により $V$ は正則である。
$Y^\nu$ は Noether なので
$Y^\nu \setminus V = \{y_1, \ldots, y_m\}$ は有限である。
各 $i$ について $\mathcal{O}_{Y^\nu, y_i}^h$ を Hensel 局所環とする。
このとき $X \times_Y \Spec(\mathcal{O}_{Y^\nu, y_i}^h)$ は
$\Spec(\mathcal{O}_{Y^\nu, y_i}^h)$ の正則 alteration である
（いくつかの詳細は省略する）。Lemma
\ref{spaces-resolve-lemma-regular-alteration-implies-local} により
$\mathcal{O}_{Y^\nu, y_i}^h$ の完備化は正規である。
従って (1) $\Rightarrow$ (4) を得る。

\medskip\noindent
(4) を仮定し，(3) を証明する。直ちに $Y$ をその正規化で置き換えてよい。
$y_1, \ldots, y_m \in |Y|$ を特異点とする。Section
\ref{spaces-resolve-section-strategy} の通り初等 \'etale 近傍
$(V_i, v_i) \to (Y, y_i)$ の族を選ぶ。各 $i$ について Hensel 局所環
$\mathcal{O}_{Y^\nu, y_i}^h$ は $\mathcal{O}_{V_i, v_i}$ の Hensel 化である。
したがってこれらの環の完備化は同型である。スキームに対する結果
（Resolution of Surfaces, Theorem \ref{resolve-theorem-resolve}）により，
$v_i$ の上の点のみをブローアップする正規化ブローアップの有限列
$$
X_{i, n_i} \to X_{i, n_i - 1} \to \ldots \to V_i
$$
で，$X_{i, n_i}$ が正則となるものが存在する。Lemma
\ref{spaces-resolve-lemma-equivalence-sequence-normalized-blowups} により，
正規化ブローアップの列
$$
X_n \to X_{n - 1} \to \ldots \to X_1 \to Y
$$
が存在し，もちろん $X_n$ も正則である（局所環を見ればよい）。
これで証明が完了する。
\end{proof}









\section{例}
\label{spaces-resolve-section-examples}

\noindent
本章で先に得た結果に関連するいくつかの例を挙げる。

\begin{example}
\label{spaces-resolve-example-factorial}
\begin{reference}
\cite[4(c)]{Samuel-UFD}
\end{reference}
$k$ を体とする。$r, s, t$ が互いに素な整数ならば，環
$A = k[x, y, z]/(x^r + y^s + z^t)$ は UFD である。
実際，$x^r + y^s + z^t$ は既約なので $A$ は整域である。
元 $z$ は素元，すなわち $A$ の素イデアルを生成する。
一方，ある $e$ に対して $t = 1 + ers$ ならば
$$
A[1/z] \cong k[x', y', 1/z]
$$
である。ただし $x' = x/z^{es}$，$y' = y/z^{er}$，および
$z = (x')^r + (y')^s$ である。従って $A[1/z]$ は多項式環の局所化であり，
ゆえに UFD である。Nagata の議論から $A$ も UFD である。
Algebra, Lemma \ref{algebra-lemma-invert-prime-elements} を参照せよ。
$t$ が $rs$ を法として $1$ と合同でない場合にも，同様の議論を与えられる。
\end{example}

\begin{example}
\label{spaces-resolve-example-completion-not-factorial}
\begin{reference}
一般の場合の局所 Picard 群の非消滅については
\cite{Brieskorn} および \cite{Lipman-rational} を参照せよ。
\end{reference}
$1 < r < s < t$ が互いに素な整数で $2, 3, 5$ ではないとき，環
$A = \mathbf{C}[[x, y, z]]/(x^r + y^s + z^t)$ は UFD ではない。
例えば特別な場合
$A = \mathbf{C}[[x, y, z]]/(x^2 + y^5 + z^7)$ を考える。
写像
$$
\psi_\zeta : \mathbf{C}[[x, y, z]]/(x^2 + y^5 + z^7) \to \mathbf{C}[[t]]
$$
を
$$
x \mapsto t^7,\quad
y \mapsto t^3,\quad
z \mapsto -\zeta t^2(1 + t)^{1/7}
$$
で定める。ただし $\zeta$ は $7$ 乗根である。$\psi_\zeta$ の核
$\mathfrak p_\zeta$ は高さ 1 の素イデアルなので，$A$ が UFD ならば
単項であり，ある $f_\zeta \in \mathbf{C}[[x, y, z]]$ で生成される。
$V(x^3 - y^7) = \bigcup V(\mathfrak p_\zeta)$ であり，
$A/(x^3 - y^7)$ は閉点の外で被約であることに注意せよ。
従って，なお $A$ が UFD であると仮定すれば
$$
\prod\nolimits_\zeta f_\zeta = u(x^3 - y^7) + a(x^2 + y^5 + z^7)
\quad\text{in}\quad
\mathbf{C}[[x, y, z]]
$$
を得る。ただし $u \in \mathbf{C}[[x, y, z]]$ は単元，
$a \in \mathbf{C}[[x, y, z]]$ である。定数倍して
$u(0, 0, 0) = 1$ と仮定してよい。左辺は位数 $7$ で消えることに注意せよ。
従って $a = - x \bmod \mathfrak m^2$ である。しかしこのとき右辺には
項 $xy^5$ が現れるが，左辺には現れない。これは矛盾である。
\end{example}

\begin{example}
\label{spaces-resolve-example-not-blow-up}
スキームではない修正 $X \to S = \Spec(A)$ をもつ優秀な $2$ 次元
Noether 局所環が存在する。構成を概説する。$X$ を
$\mathbf{C}$ 上の正規曲面で，特異点 $x \in X$ をただ一つもつものとする。
例外ファイバー $C = \pi^{-1}(x)_{red}$ が滑らかな射影曲線となるような
特異点解消 $\pi : X' \to X$ が存在すると仮定する。さらに，
$\mathcal{O}_C(nc)$ が $\Pic(X') \to \Pic(C)$ の像に属するならば
$n = 0$ となるような点 $c \in C$ が存在すると仮定する。
$X'' \to X'$ を非特異点 $c$ におけるブローアップとする。
$C' \subset X''$ を $C$ の狭義変換とし，$E \subset X''$ を例外ファイバーとする。
Artin の結果（\cite{ArtinII}；例えば \cite{Mumford-topology} を用いれば
$C'$ の法束が負であることが分かる）により，$X''$ の曲線 $C'$ を収縮して
代数空間 $X'''$ を得ることができる。図式は
$$
\xymatrix{
& X'' \ar[ld] \ar[rd] \\
X' \ar[rd] &  & X''' \ar[ld] \\
& X
}
$$
である。$X'''$ はスキームではないと主張する。これにより所望の例が得られる。
実際，$X'''$ がスキームであることと，$A = \mathcal{O}_{X, x}$ への
$X'''$ の基底変換がスキームであることは同値である（詳細は省略する）。
% P10-E0416: frozen source writes "where" instead of "were" in this counterfactual.
もし $X'''$ がスキームであるならば，$X'''$ における $C'$ の像は
アフィン近傍をもつ。この近傍の補集合は $X'''$ 上の有効 Cartier 因子となる
（$X'''$ は $1$ 点を除いて非特異だからである）。この有効 Cartier 因子は，
$E$ と交わり $C'$ を避ける $X''$ 上の有効 Cartier 因子に対応する。
$X'$ における像を取れば，$C$ と（集合論的に）$c$ で交わる有効 Cartier
因子を得る。仮定により $c$ のいかなる倍数も Cartier 因子の制限ではないので，
これは不可能である。

\medskip\noindent
最後に，このような特異曲面 $X$ を見つけなければならない。
$\mathbf{A}^3_\mathbf{C} = \Spec(\mathbf{C}[x, y, z])$ において
$$
x^3 + y^3 + z^3 + x^4 + y^4 + z^4 = 0
$$
で定まるアフィン曲面を $X$ とし，特異点を $(0, 0, 0)$ とすればよい。
このとき $(0, 0, 0)$ は唯一の特異点である。$(0, 0, 0)$ に対応する
極大イデアルで $X$ をブローアップすると，それぞれ滑らかなアフィン曲面
$$
1 + s^3 + t^3 + x(1 + s^4 + t^4) = 0
$$
と同型な三つのチャートを得る。この曲面は非特異で，例外因子 $C$ は
$x = 0$ で与えられる。読者は $C$ が楕円曲線であることに気付くだろう。
さらに，$(0, 0, 0)$ からの射影により $X$ は有理的であることが分かる。
または，ブローアップの方程式で $x$ について解いてもよい。
非特異有理曲面の Picard 群は可算である一方，複素数体上の楕円曲線の
Picard 群は非可算である。従って，述べた性質をもつ閉点 $c$ を取れる。
\end{example}
